\subsection{缺点：光学参数互补与外部标定缺失}
\label{sec:limitations}

连续波段重估后的厚度仍然分散，厚度重估范围也随光学假设扩展。这使评价重点从曲线拟合好坏转向厚度是否被观测唯一确定。

碳化硅补充五折连续留段检验后，厚度分布于$5.53$--$13.51\,\mu\mathrm{m}$，极差与中位数之比为$0.7604$。部分分支超出原范围，因此将全部有效切分纳入范围；图\ref{fig:q2-thickness}保留了这种波段依赖。
% src:求解/问题2/结果/验证.json:五折跨切分稳定性.最小厚度_微米;五折跨切分稳定性.最大厚度_微米;五折跨切分稳定性.相对极差;五折跨切分稳定性.有效折数
% src:交接/实验记录.json:505.尝试;505.现象;505.决定

波段依赖反映厚度与未知折射率变化的互相补偿，即不同厚度也能配出相近条纹。纳入损失接近最小值的近优候选及参数扰动后，范围扩展至$0.64$--$18.11\,\mu\mathrm{m}$；这一跨度由已检验情景的最外端点确定，反映现有观测对厚度的约束较弱。
% src:求解/问题2/结果/厚度结果.json:条件区间_微米;条件区间组成

物理解释还受观测条件制约：相干性、空间重叠、仪器分辨率和偏振状态四类资料均缺失（表\ref{tab:07}）。完整往返场模型在部分区段改善谱形，总体误差却高于两束；物理可观测条件仍有缺项，碳化硅沿用问题二基准（图\ref{fig:20}）。
% src:求解/问题3/结果/条件判定.json:必要条件;碳化硅条件分支.触发
% src:交接/数据档案.json:未随附件提供的建模输入

区间覆盖也存在缺口：问题二、三的反射率测试经验覆盖率分别为$49.06\%$、$73.91\%$，前者低于$75\%$校准分位水平，后者低于$90\%$名义目标（图\ref{fig:q2-coverage}、表\ref{tab:09}）。校准残差给出的区间宽度未能充分容纳测试偏差，而真实附件的厚度真值也缺少独立参照。推广应从光学常数、标准厚度样片和重复测量补足这些缺口。
% src:求解/问题2/结果/区间覆盖.json:测试经验覆盖率;校准经验分位点
% src:求解/问题3/结果/回炉轮3候选/仲裁闭环/20260910_131501_81623_建模公式/厚度结果.json:预测区间.4.经验覆盖率;预测区间.4.平均区间宽度_比例;预测区间.5.经验覆盖率;预测区间.5.平均区间宽度_比例;预测区间.6.经验覆盖率;预测区间.6.平均区间宽度_比例;预测区间.7.经验覆盖率;预测区间.7.平均区间宽度_比例;预测区间.20.经验覆盖率;预测区间.20.平均区间宽度_比例;预测区间.21.经验覆盖率;预测区间.21.平均区间宽度_比例;预测区间.22.经验覆盖率;预测区间.22.平均区间宽度_比例;预测区间.23.经验覆盖率;预测区间.23.平均区间宽度_比例;预测区间.36.经验覆盖率;预测区间.36.平均区间宽度_比例;预测区间.37.经验覆盖率;预测区间.37.平均区间宽度_比例;预测区间.38.经验覆盖率;预测区间.38.平均区间宽度_比例;预测区间.39.经验覆盖率;预测区间.39.平均区间宽度_比例;预测区间.52.经验覆盖率;预测区间.52.平均区间宽度_比例;预测区间.53.经验覆盖率;预测区间.53.平均区间宽度_比例;预测区间.54.经验覆盖率;预测区间.54.平均区间宽度_比例;预测区间.55.经验覆盖率;预测区间.55.平均区间宽度_比例;预测区间.4.名义覆盖率
% src:求解/问题3/结果/回炉轮3候选/仲裁闭环/20260910_131501_81623_建模公式/厚度结果.json:逐块评分.4.测试点数;逐块评分.5.测试点数;逐块评分.6.测试点数;逐块评分.7.测试点数;逐块评分.20.测试点数;逐块评分.21.测试点数;逐块评分.22.测试点数;逐块评分.23.测试点数;逐块评分.36.测试点数;逐块评分.37.测试点数;逐块评分.38.测试点数;逐块评分.39.测试点数;逐块评分.52.测试点数;逐块评分.53.测试点数;逐块评分.54.测试点数;逐块评分.55.测试点数
% src:交接/数据档案.json:未随附件提供的建模输入
